\scanpage{154}
$K$ corps CM \inserted{galoisien} sur $\mathbb Q$ groupe de Galois $G$,
$\unsure{\gamma|\sigma}=\{1,\rho\}\simeq\mathbb Z/2\mathbb Z\subset G$.
\margin{$\rho$ la conj. complexe}
$\overline{\mathbb Q}$ \inserted{une} clôture algébrique de $\mathbb Q$
\deleted{\unsure{contenant} $K$}.

\deleted{$I$ ensemble des plongts de $K$ dans $\overline{\mathbb Q}$.
$G$ y opère par transport de structure $g.\sigma=\sigma\circ g^{-1}$
de façon simplement transitive.}

$M$ groupe des fonctions \deleted{\unsure{\varphi}} $\varphi:\unsure{G|I}-\mathbb Z$ telles que
\[
\varphi(\rho\sigma)=-\varphi(\sigma)\qquad\text{pour }\forall\,\sigma\in G
\]
C'est un $G$-module, qui définit \inserted{par descente et dualité} un tore $T$ sur $\mathbb Q$.
\margin{\uline{NB} $M\simeq\check M$ comme $G$-modules.}

\textcircled{1} On veut calculer le sous-groupe \inserted{$H^1(\mathbb Q,T)'$} de
$H^1(\mathbb Q,T)$ formé des classes qui sont nulles en toute \uline{place finie}.
On \unsure{trouve} par \deleted{si \unsure{$I$} $P=I/\gamma$, alors l'on a}
une \deleted{isom} \inserted{suite exacte} canonique
\diagram{\begin{tikzcd}[column sep=large]
0\arrow[r]&\unsure{\widehat H^1|H^1}(\mathbb Q,T)'\arrow[r,"\alpha"]&
\widehat H_0(G,Y_\infty\otimes_{\mathbb Z}\check M)\arrow[r,"\beta"]&
\widehat H_0(G,\check M)\arrow[r]&0
\end{tikzcd}}
où $Y_\infty$ est le groupe libre engendré par les places à l'$\infty$ de $K$,
et où $\beta$ est déduit de l'augmentation naturelle $Y_\infty\longrightarrow\mathbb Z$
(somme des \unsure{éléments|entiers}).

Si on prend pour $\overline{\mathbb Q}$ la clôture algébrique de $\mathbb Q$ dans $\mathbb C$,
on a $M\simeq\mathbb Z'\otimes_{\gamma}G$, où $\mathbb Z'$ est le $\gamma$-module
$\mathbb Z$ avec $\rho$ opérant par $-\operatorname{id}_{\mathbb Z}$. On a ainsi
\diagram{\begin{tikzcd}[column sep=large]
0\arrow[r,no head]&H^1(\mathbb Q,T)'\arrow[r,no head]&
\widehat H_0(\gamma,Y'_\infty)\arrow[r]\arrow[d,phantom,"\wr" description]&
\widehat H_0(\gamma,\mathbb Z')\arrow[r]\arrow[d,equal]&0\\
&&Y_\infty\otimes\mathbb Z/2\mathbb Z\arrow[r]&\mathbb Z/2\mathbb Z&\deleted{\unsure{surj.}}
\end{tikzcd}}
et donc on trouve
\[
H^1(\mathbb Q,T)'\simeq\text{gpe des combinaisons linéaires
\inserted{à coeff. dans $\mathbb F_2$} de places à l'infini de $K$, à somme des coefficients nulle.}
\]
\margin{153}
